\documentclass[11pt]{article}

\usepackage[margin=1in]{geometry}
\usepackage{amsmath,amssymb,amsthm}
\usepackage{mathtools}
\usepackage[T1]{fontenc}
\usepackage{lmodern}
\usepackage{enumitem}
\usepackage{microtype}

\allowdisplaybreaks
\setlength{\parindent}{0pt}
\setlength{\parskip}{5pt}

\newcommand{\pos}{\operatorname{pos}}
\newcommand{\U}{\mathsf{U}}
\newcommand{\D}{\mathsf{D}}

\begin{document}

\begin{center}
    {\Large\bfseries MATH 465 -- Homework 5 Solutions}\\[4pt]
    Due October 8, 2026\\[8pt]
    \textbf{Name:} \underline{\hspace{2.5in}}
\end{center}

\textbf{Acknowledgment.} I consulted the homework prompt and the course
lecture notes. No students were consulted. Generative AI assistance was used
in preparing this write-up.

\section*{1. Finite differences and interpolation}

Define
\[
    q(t)=p(2t).
\]
Then \(q\) has degree at most \(4\), and the given values of \(p\) give the
following forward-difference table:
\[
\begin{array}{c|ccccc}
t & 0&1&2&3&4\\ \hline
q(t) & 1&4&13&40&109\\
\Delta q(t) & 3&9&27&69\\
\Delta^2q(t) & 6&18&42\\
\Delta^3q(t) & 12&24\\
\Delta^4q(t) & 12
\end{array}
\]
The fourth difference is constant because \(q\) has degree at most \(4\).

\subsection*{(a) Values at \(10\) and \(-2\)}

To extend the table to \(t=5\), first extend the successive differences:
\[
\begin{aligned}
\Delta^3q(4)&=\Delta^3q(3)+\Delta^4q(3)=24+12=36,\\
\Delta^2q(4)&=\Delta^2q(3)+\Delta^3q(4)=42+36=78,\\
\Delta q(4)&=\Delta q(3)+\Delta^2q(4)=69+78=147.
\end{aligned}
\]
Therefore
\[
    q(5)=q(4)+\Delta q(4)=109+147=256,
\]
and hence
\[
    \boxed{p(10)=256}.
\]

Newton's forward interpolation formula gives
\[
    q(t)=q(0)+\Delta q(0)\binom{t}{1}
    +\Delta^2q(0)\binom{t}{2}
    +\Delta^3q(0)\binom{t}{3}
    +\Delta^4q(0)\binom{t}{4},
\]
so in this case
\[
    q(t)=1+3\binom{t}{1}+6\binom{t}{2}
             +12\binom{t}{3}+12\binom{t}{4}.
\]
At \(t=-1\), the generalized binomial identity
\(\binom{-1}{k}=(-1)^k\) gives
\[
    q(-1)=1-3+6-12+12=4.
\]
Since \(p(-2)=q(-1)\),
\[
    \boxed{p(-2)=4}.
\]

\subsection*{(b) Formula and uniqueness}

Substituting \(t=x/2\) into the interpolation formula yields
\[
\boxed{
p(x)=1+3\binom{x/2}{1}+6\binom{x/2}{2}
        +12\binom{x/2}{3}+12\binom{x/2}{4}.
}
\]
This polynomial is uniquely determined by the five data values. Indeed, if
\(p_1\) and \(p_2\) both have degree at most \(4\) and satisfy the given
conditions, then \(p_1-p_2\) has the five distinct roots
\(0,2,4,6,8\). A nonzero polynomial of degree at most \(4\) cannot have five
distinct roots, so \(p_1-p_2\) is identically zero.

\section*{2. Integer-valued polynomials}

\subsection*{(a) Integer values at every integer}

Let \(\Delta f(x)=f(x+1)-f(x)\). The Newton expansion for a polynomial of
degree at most \(d\) is
\[
    p(x)=\sum_{k=0}^{d}\Delta^k p(0)\binom{x}{k},
    \qquad
    \binom{x}{k}=\frac{x(x-1)\cdots(x-k+1)}{k!}.
\]
To justify this expansion, note that
\(\Delta\binom{x}{k}=\binom{x}{k-1}\), with
\(\binom{x}{0}=1\), so the successive differences of the right-hand side
at \(x=0\) are exactly \(\Delta^k p(0)\). The two sides are polynomials of
degree at most \(d\) with the same values at \(0,1,\ldots,d\), and therefore
they are equal.

Because \(p(0),\ldots,p(d)\) are integers, every
\(\Delta^k p(0)\) is an integer: each finite difference is an integer linear
combination of these given values. For a nonnegative integer \(m\),
\(\binom{m}{k}\) is an integer in the usual sense. For a negative integer
\(m\),
\[
    \binom{m}{k}
    =(-1)^k\binom{-m+k-1}{k},
\]
which is also an integer. Thus every term in the Newton expansion is an
integer when \(x=m\in\mathbb Z\), and consequently
\[
    \boxed{p(m)\in\mathbb Z\quad\text{for every }m\in\mathbb Z.}
\]

\subsection*{(b) Ordinary coefficients need not be integers}

No. For example, take \(d=2\) and
\[
    p(x)=\binom{x}{2}=\frac{x^2-x}{2}.
\]
Then
\[
    p(0)=0,\qquad p(1)=0,\qquad p(2)=1
\]
are integers, and part (a) shows that \(p(m)\) is an integer for every
integer \(m\). However, the coefficients of \(x^2\) and \(x\) in the
ordinary power basis are \(1/2\) and \(-1/2\), respectively. Hence
integer-valuedness does not imply integer coefficients in the power basis.

\section*{3. A Stirling-number summation identity}

\subsection*{(a) Proof of the identity}

Using the permitted identity
\[
    j^m=\sum_{k=0}^{m}S(m,k)(j)_k,
    \qquad
    (j)_k=j(j-1)\cdots(j-k+1),
\]
we obtain
\[
\begin{aligned}
\sum_{j=0}^{N-1}j^m
  &=\sum_{j=0}^{N-1}\sum_{k=0}^{m}S(m,k)(j)_k\\
  &=\sum_{k=0}^{m}S(m,k)\sum_{j=0}^{N-1}k!\binom{j}{k}\\
  &=\sum_{k=0}^{m}k!S(m,k)\sum_{j=0}^{N-1}\binom{j}{k}.
\end{aligned}
\]
The hockey-stick identity gives
\[
    \sum_{j=0}^{N-1}\binom{j}{k}=\binom{N}{k+1}.
\]
This remains true for \(k=0\), and when \(N=0\) both sides of the desired
identity are empty sums or zero binomial coefficients. Therefore
\[
\boxed{
\sum_{j=0}^{N-1}j^m
    =\sum_{k=0}^{m}k!S(m,k)\binom{N}{k+1}.
}
\]

\subsection*{(b) The fourth-power sum}

The relevant Stirling numbers of the second kind are
\[
    S(4,0)=0,\qquad S(4,1)=1,\qquad S(4,2)=7,\qquad
    S(4,3)=6,\qquad S(4,4)=1.
\]
Multiplying by the corresponding factorials gives
\[
\begin{array}{c|ccccc}
k&0&1&2&3&4\\ \hline
k!S(4,k)&0&1&14&36&24
\end{array}
\]
and hence
\[
\boxed{
\sum_{j=0}^{N-1}j^4
   =\binom{N}{2}+14\binom{N}{3}
      +36\binom{N}{4}+24\binom{N}{5}.
}
\]

\section*{4. A nonhomogeneous linear recurrence}

\subsection*{(a) A resonant trial sequence cannot work}

If \(h_n=A2^n\), then the left-hand side of the recurrence is
\[
\begin{aligned}
h_n-3h_{n-1}+2h_{n-2}
  &=A2^{n-2}(4-6+2)\\
  &=0.
\end{aligned}
\]
The right-hand side is \(2^n\ne0\). Thus no sequence of the form
\(h_n=A2^n\), with \(A\) constant, satisfies the recurrence.

\subsection*{(b) Closed formula}

The homogeneous recurrence has characteristic polynomial
\[
    \lambda^2-3\lambda+2=(\lambda-1)(\lambda-2),
\]
so
\[
    h_n^{(h)}=C+D2^n.
\]
Because \(2\) is a root of the characteristic polynomial, use the resonant
particular trial solution \(h_n^{(p)}=E\,n2^n\). Substitution gives
\[
\begin{aligned}
h_n^{(p)}-3h_{n-1}^{(p)}+2h_{n-2}^{(p)}
 &=E2^{n-2}\bigl(4n-6(n-1)+2(n-2)\bigr)\\
 &=E2^{n-1}.
\end{aligned}
\]
Equating this with \(2^n\) gives \(E=2\). Therefore
\[
    h_n=C+D2^n+2n2^n.
\]
The initial conditions imply
\[
    C+D=0,\qquad C+2D+4=1.
\]
Thus \(D=-3\) and \(C=3\), and
\[
\boxed{
    h_n=3+(2n-3)2^n\qquad(n\geq0).
}
\]
The formula gives \(h_0=0\) and \(h_1=1\), and the substitution above shows
that it satisfies the recurrence.

\section*{5. Paths bounded below by a horizontal line}

Write \(\U=(1,1)\) and \(\D=(1,-1)\). Let \(A_r\) denote the number of paths
from \((0,0)\) to \((2n,0)\) that stay in \(y\geq-r\).

\subsection*{(a) Paths staying in \(y\geq-r\)}

There are \(\binom{2n}{n}\) unrestricted paths, since a path ending at
height \(0\) has \(n\) up-steps and \(n\) down-steps. A path that violates
the condition \(y\geq-r\) must first hit the line \(y=-(r+1)\). Reflect the
portion of the path from its beginning through this first hitting point
across the line \(y=-(r+1)\). The starting point \((0,0)\) is sent to
\((0,-2r-2)\), while the endpoint remains \((2n,0)\).

This reflection is a bijection between bad paths and unrestricted paths from
\((0,-2r-2)\) to \((2n,0)\). Such a path has \(n+r+1\) up-steps and
\(n-r-1\) down-steps, so there are
\(\binom{2n}{n+r+1}\) of them. Consequently,
\[
\boxed{
    A_r=\binom{2n}{n}-\binom{2n}{n+r+1}.
}
\]
The stated binomial convention makes this valid also when \(r=n\).

\subsection*{(b) Paths with minimum height exactly \(-r\)}

For \(r\geq1\), the paths staying above \(-r\) but not above \(-(r-1)\)
are counted by \(A_r-A_{r-1}\). Therefore
\[
\begin{aligned}
A_r-A_{r-1}
 &=\left[\binom{2n}{n}-\binom{2n}{n+r+1}\right]
   -\left[\binom{2n}{n}-\binom{2n}{n+r}\right]\\
 &=\binom{2n}{n+r}-\binom{2n}{n+r+1}.
\end{aligned}
\]
When \(r=0\), every path with minimum height exactly \(0\) is simply a path
staying in \(y\geq0\), and the same expression is \(A_0\). Hence for every
\(0\leq r\leq n\),
\[
\boxed{
\#\{\text{paths with minimum height exactly }-r\}
 =\binom{2n}{n+r}-\binom{2n}{n+r+1}.
}
\]

\section*{6. Triangulations and Catalan numbers}

\subsection*{(a) The recurrence}

Consider the triangle incident to the boundary edge \(v_0v_{n+1}\). Its
third vertex is \(v_{k+1}\) for a unique \(k\) with \(0\leq k\leq n-1\).
The diagonal sides of this triangle split the polygon into two independent
polygons:
\[
    v_0,v_1,\ldots,v_{k+1}
    \quad\text{and}\quad
    v_{k+1},v_{k+2},\ldots,v_{n+1}.
\]
The first has \(k+2\) vertices and hence \(t_k\) triangulations; the second
has \(n-k+1\) vertices and hence \(t_{n-k-1}\) triangulations. Conversely,
choosing triangulations of these two subpolygons and adjoining the
distinguished triangle gives a unique triangulation of the original
polygon. Thus
\[
\boxed{
    t_0=1,\qquad
    t_n=\sum_{k=0}^{n-1}t_k\,t_{n-1-k}\quad(n\geq1).
}
\]

The Catalan numbers are characterized by the same recurrence
\[
    C_0=1,\qquad C_n=\sum_{k=0}^{n-1}C_kC_{n-1-k}.
\]
Induction on \(n\) therefore gives \(t_n=C_n\) for every \(n\). Equivalently,
the generating function \(C(z)=\sum_{n\geq0}C_nz^n\) satisfies
\(C(z)=1+zC(z)^2\), whose solution with constant term \(1\) is
\[
    C(z)=\frac{1-\sqrt{1-4z}}{2z}
       =\sum_{n\geq0}\frac{1}{n+1}\binom{2n}{n}z^n.
\]
Thus \(t_n=C_n=\frac{1}{n+1}\binom{2n}{n}\).

\subsection*{(b) A prescribed diagonal}

For \(1\leq r\leq n-1\), the diagonal \(v_0v_{r+1}\) splits the polygon into
the subpolygons
\[
    v_0,v_1,\ldots,v_{r+1}
    \quad\text{and}\quad
    v_0,v_{r+1},v_{r+2},\ldots,v_{n+1}.
\]
Their triangulations can be chosen independently. They have \(r+2\) and
\(n-r+2\) vertices, respectively, so the number of triangulations containing
the diagonal is
\[
\boxed{t_rt_{n-r}=C_rC_{n-r}.}
\]

\section*{7. Fully bracketed expressions and ballot sequences}

\subsection*{(a) The first construction is not injective}

The two distinct bracketings
\[
    ((a_1a_2)a_3)
    \qquad\text{and}\qquad
    (a_1(a_2a_3))
\]
both produce the word \(++--\): in either case there are two opening
parentheses followed, after deleting the letters, by two closing parentheses. Hence the construction in part (a) is not injective.

\subsection*{(b) The second construction gives a ballot sequence}

There are \(n\) opening parentheses and \(n\) closing parentheses in a
bracketing on \(n+1\) letters. After deleting all closing parentheses and
the final letter \(a_{n+1}\), the construction therefore produces a word of
length \(2n\), containing \(n\) plus signs and \(n\) minus signs.

To check the prefix condition, view a bracketing as an ordered full binary
tree: each opening parenthesis is an internal node and each letter is a
leaf. Reading the expression from left to right is preorder traversal. Start
with one available child slot. Reading an opening parenthesis fills one slot
and creates two, so the number of available slots increases by one. Reading a
letter fills one slot, so the number decreases by one. Before the final leaf
\(a_{n+1}\) is read, at least one slot remains. Thus every prefix of the
constructed word satisfies
\[
    1+\#(+)-\#(-)\geq1,
\]
or equivalently \(\#(+)\geq\#(-)\). The word is therefore a ballot sequence.

\subsection*{(c) The inverse and the enumeration}

Let \(w\) be a ballot sequence of length \(2n\) with \(n\) plus signs and
\(n\) minus signs. Append one additional minus sign, which will represent
the final letter \(a_{n+1}\). Reconstruct an ordered full binary tree by
processing this extended word from left to right:
\begin{itemize}[leftmargin=2em]
    \item on a \(+\), fill the leftmost available slot with an internal node
    and create its two ordered child slots;
    \item on a \(-\), fill the leftmost available slot with a leaf.
\end{itemize}
After any prefix of \(w\), the number of slots is
\(1+\#(+)-\#(-)\geq1\); after the final appended minus it is zero. Thus the
algorithm is well-defined and produces a full binary tree with \(n\) internal
nodes and \(n+1\) leaves. Label the leaves from left to right by
\(a_1,\ldots,a_{n+1}\), and write an internal node with left and right
subexpressions \(U,V\) as \((UV)\). This yields a fully bracketed expression.

For an expression, the same slot procedure reads its preorder tree and
recovers exactly the original tree. Conversely, the reconstructed tree has
the prescribed preorder word, so deleting closing parentheses and the final
leaf returns \(w\). The two constructions are therefore inverse bijections.

The number of ballot sequences of length \(2n\) is, by the reflection
principle,
\[
    \binom{2n}{n}-\binom{2n}{n+1}
    =\frac{1}{n+1}\binom{2n}{n}=C_n.
\]
Consequently, the number of fully bracketed expressions on \(n+1\) letters is
\[
\boxed{C_n=\frac{1}{n+1}\binom{2n}{n}.}
\]

\section*{8. Dyck paths counted by peaks}

\subsection*{(a) The first three polynomials}

For semilength \(1\), the only path is \(\U\D\), with one peak:
\[
    \boxed{P_1(q)=q}.
\]
For semilength \(2\), the paths are \(\U\U\D\D\) and
\(\U\D\U\D\), with one and two peaks:
\[
    \boxed{P_2(q)=q+q^2}.
\]
For semilength \(3\), the five paths have peak counts
\(1,2,2,2,3\), respectively. Thus
\[
    \boxed{P_3(q)=q+3q^2+q^3}.
\]

\subsection*{(b) A first-return recurrence}

Take a Dyck path \(D\) of semilength \(n+1\), and write its unique
first-return decomposition as
\[
    D=\U A\D B,
\]
where \(A\) is a Dyck path of semilength \(k\) and \(B\) is a Dyck path of
semilength \(n-k\), for some \(0\leq k\leq n\).

If \(k=0\), the initial \(\U\D\) is a peak and contributes a factor \(q\);
the remaining path \(B\) contributes \(P_n(q)\). This case contributes
\(qP_n(q)\). If \(k\geq1\), then \(A\) begins with an up-step, so the initial
\(\U\) is not immediately followed by a down-step. The peaks are exactly
those in \(A\) and \(B\), giving \(P_k(q)P_{n-k}(q)\). Summing over all
possibilities,
\[
\boxed{
    P_{n+1}(q)=qP_n(q)+\sum_{k=1}^{n}P_k(q)P_{n-k}(q)
    \qquad(n\geq0).
}
\]

\subsection*{(c) Peaks and leaves in an ordered rooted tree}

For an ordered rooted tree with \(n+1\) vertices, traverse each edge away from
the root and record an up-step when an edge is traversed downward and a
down-step when that edge is traversed upward. Each of the \(n\) edges is
traversed twice, so this gives a Dyck path of length \(2n\).

An up-step is immediately followed by a down-step exactly when the traversal
reaches a vertex and has no child to visit before returning. This occurs
exactly at a leaf. Conversely, every nonroot leaf creates such a peak.
Since \(n\geq1\), the root is not a leaf, so
\[
\boxed{\text{the number of peaks equals the number of leaves of the tree}.}
\]

\section*{9. 132-avoiding permutations and Rothe diagrams}

For a permutation \(\sigma\), write
\[
    \pos(j)=\sigma^{-1}(j).
\]
Thus a cell in row \(i\) and column \(j\) of the Rothe diagram is present
exactly when
\[
    j<\sigma(i)\quad\text{and}\quad \pos(j)>i.
\]
Let \(d_i\) be the number of cells in row \(i\).

\subsection*{(a) Characterization by row lengths}

\paragraph{132-avoidance implies left-justified rows.}
Suppose \(\sigma\) is 132-avoiding, and suppose row \(i\) contains a cell in
column \(j\). If some \(\ell<j\) were not in row \(i\), then
\(\ell<j<\sigma(i)\), and \(\pos(\ell)<i\), while \(\pos(j)>i\). Set
\(h=\pos(\ell)\) and \(k=\pos(j)\). Then
\[
    h<i<k,\qquad \sigma(h)=\ell<j=\sigma(k)<\sigma(i),
\]
which is a 132 pattern. Therefore every row consists of its leftmost
\(d_i\) cells.

\paragraph{132-avoidance implies weakly decreasing row lengths.}
Assume the rows are left-justified. If \(d_{i+1}>d_i\), let
\(j=d_i+1\). Then \((i+1,j)\) is a cell, so
\[
    j<\sigma(i+1),\qquad \pos(j)>i+1.
\]
The cell \((i,j)\) is absent, and \(\pos(j)>i\), so \(j\geq\sigma(i)\).
Equality would force \(\pos(j)=i\), impossible; hence
\[
    \sigma(i)<j<\sigma(i+1).
\]
Using \(k=\pos(j)\), we have
\[
    i<i+1<k,\qquad \sigma(i)<j=\sigma(k)<\sigma(i+1),
\]
again a 132 pattern. Hence \(d_i\geq d_{i+1}\) for every \(i\).

\paragraph{The converse.}
Now suppose every row is left-justified and
\[
    d_1\geq d_2\geq\cdots\geq d_n\geq0.
\]
If \(i<j<k\) formed a 132 pattern, put \(c=\sigma(k)\). Since
\(c<\sigma(j)\) and \(\pos(c)=k>j\), the cell \((j,c)\) lies in row \(j\).
Left-justification gives \(d_j\geq c\), and weak decrease gives
\(d_i\geq d_j\geq c\). Thus \((i,c)\) would be a cell, which would imply
\(c<\sigma(i)\), contradicting the assumed inequality
\(\sigma(i)<c\). Therefore \(\sigma\) is 132-avoiding.

\paragraph{The bound on row lengths.}
Every cell in row \(i\) has its column value in a position strictly below
\(i\). There are only \(n-i\) such positions, and distinct cells have
distinct column values. Consequently
\[
\boxed{d_i\leq n-i\qquad(1\leq i\leq n).}
\]

\subsection*{(b) A bijection with Dyck paths}

The characterization in part (a) shows that a 132-avoiding permutation is
encoded by a sequence
\[
    d_1\geq d_2\geq\cdots\geq d_n\geq0,
    \qquad d_i\leq n-i.
\]
We first explain how to recover the permutation from such a sequence. Process
the rows from \(i=1\) to \(n\). At step \(i\), let \(\sigma(i)\) be the
smallest unused value that is strictly larger than \(d_i\):
\[
\boxed{\sigma(i)=\min\{j>d_i:\ j\text{ has not been used in rows }1,\ldots,i-1\}.}
\]
This set is nonempty. Indeed, all values at most \(d_i\) are still unused,
because earlier rows have lengths at least \(d_i\) and therefore chose values
larger than \(d_i\). The number of unused values larger than \(d_i\) is at
least
\[
    n-d_i-(i-1)\geq1
\]
by \(d_i\leq n-i\).

The values \(1,\ldots,d_i\) are then placed below row \(i\), so they give
exactly \(d_i\) cells. Any value \(j\) with \(d_i<j<\sigma(i)\) was already
used, by the minimality in the displayed rule, so it is not below row \(i\).
Thus the resulting permutation has precisely the prescribed Rothe diagram.
The same minimality condition shows that no other permutation can have these
row lengths, so the reconstruction is unique.

It remains to convert the row-length sequence into a Dyck path. Define
\[
    e_i=d_{n-i+1}\qquad(1\leq i\leq n).
\]
Then
\[
    0\leq e_1\leq e_2\leq\cdots\leq e_n,
    \qquad e_i\leq i-1.
\]
Construct the word
\[
\begin{aligned}
W(d)={}&
\D^{e_1}\U
\D^{e_2-e_1}\U
\cdots
\D^{e_n-e_{n-1}}\U
\D^{n-e_n}.
\end{aligned}
\]
Here \(\U=(1,1)\) and \(\D=(1,-1)\), and a zero exponent means that the
corresponding run is empty. The word has \(n\) up-steps and \(n\) down-steps.
Before the \(i\)-th up-step there are \(e_i\leq i-1\) down-steps, so no
prefix goes below height \(0\). Hence \(W(d)\) is a Dyck path of length
\(2n\).

Conversely, given a Dyck path, let \(e_i\) be the number of down-steps before
its \(i\)-th up-step. The Dyck condition gives \(e_i\leq i-1\), and the
sequence \(e_i\) is weakly increasing. Set
\[
    d_i=e_{n-i+1}.
\]
Then \(d_1\geq\cdots\geq d_n\geq0\) and \(d_i\leq n-i\). Reconstructing
\(\sigma\) by the smallest-unused-value rule above gives a unique
132-avoiding permutation. These two procedures are inverse, so they give an
explicit bijection between 132-avoiding permutations in \(S_n\) and Dyck
paths of length \(2n\). Therefore
\[
\boxed{\#\{\text{132-avoiding permutations in }S_n\}
      =C_n=\frac{1}{n+1}\binom{2n}{n}.}
\]

\section*{10. Monomials in a product of partial sums}

\[
    Q_n(x_1,\ldots,x_n)=\prod_{i=1}^{n}(x_1+\cdots+x_i).
\]
For a monomial \(x_1^{b_1}\cdots x_n^{b_n}\), necessarily
\(\sum_{i=1}^n b_i=n\), because one variable is chosen from each of the
\(n\) factors.

\subsection*{(a) The prefix condition}

Suppose the monomial occurs. In the first \(j\) factors, every selected
variable has index at most \(j\). Those \(j\) selections therefore contribute
to the total exponent of one of \(x_1,\ldots,x_j\). Consequently,
\[
    b_1+\cdots+b_j\geq j
    \qquad(1\leq j\leq n).
\]

Conversely, assume \(\sum_{i=1}^n b_i=n\) and all these prefix inequalities
hold. Assign a variable to each factor greedily from \(i=1\) through \(n\),
using exactly \(b_j\) assignments of \(x_j\). At step \(i\), there must be an
unfilled demand among \(x_1,\ldots,x_i\). Otherwise all demands in this set
would already have been filled by the first \(i-1\) factors, which would imply
\(b_1+\cdots+b_i\leq i-1\), contradicting the prefix condition. Choose any
such \(x_j\) with \(j\leq i\) from the \(i\)-th factor. After all \(n\)
steps, every demand is filled, so the monomial occurs in \(Q_n\). Thus
\[
\boxed{
x_1^{b_1}\cdots x_n^{b_n}\text{ occurs in }Q_n
\iff
\sum_{i=1}^n b_i=n\ \text{and}\ 
\sum_{i=1}^j b_i\geq j\ (1\leq j\leq n).
}
\]

\subsection*{(b) Fixing the exponent of \(x_1\)}

Encode an exponent vector satisfying part (a) by the word
\[
    W(b)=\U^{b_1}\D\,\U^{b_2}\D\cdots
          \U^{b_n}\D.
\]
There are \(n\) up-steps in total and \(n\) down-steps. After the \(j\)-th
down-step, the height is
\[
    (b_1+\cdots+b_j)-j,
\]
so the prefix inequalities in part (a) are exactly the Dyck-path condition.
If \(b_1=r\), the path begins with exactly \(r\) up-steps followed by a
down-step. Thus we count Dyck paths of semilength \(n\) whose initial run of
up-steps has length exactly \(r\).

After the forced prefix \(\U^r\D\), the remaining path has length
\(2n-r-1\), starts at height \(r-1\), ends at height \(0\), and contains
\(n-r\) up-steps and \(n-1\) down-steps. Ignoring the nonnegativity
condition, there are
\(\binom{2n-r-1}{n-r}\) such paths. Reflecting the initial segment through
the first visit to \(y=-1\) is a bijection from the bad paths to paths
starting at height \(-r-1\) and ending at \(0\). A reflected path has \(n\)
up-steps, so the number of bad paths is
\(\binom{2n-r-1}{n}\). Therefore the number of exponent vectors, and hence of
distinct monomials, is
\[
\boxed{
\#\{\text{monomials with exponent of }x_1\text{ equal to }r\}
 =\binom{2n-r-1}{n-r}-\binom{2n-r-1}{n},
\qquad 1\leq r\leq n.
}
\]

\end{document}


